\label{chap:83-09}

% Unidad U012. Biografía estructural completa de pi.

% La biografía de pi que seguía en esta residencia es exactamente la ya
% publicada en el capítulo 8, salvo el namespace de sus etiquetas. Se conserva
% en la fuente por su procedencia, pero no se imprime dos veces.
\iffalse

\section{Construcción del carácter de clausura y trayectorias genealógicas estructural de \(\pi\)}
\label{p12-83-c9-s1:chap:biografia-pi-autonoma}

\subsection{Separación causal entre semilla, carácter y evaluación}

La construcción se divide en tres estratos:
\[
 \text{selección finita}
 \longrightarrow
 \text{construcción del carácter de clausura}
 \longrightarrow
 \text{evaluación arquimediana}.
\]
La expansión decimal no interviene en el primer estrato. La cadena de tipos
es
\begin{equation}
 104\,976
 \longrightarrow468\,\mathcal U_6
 \xrightarrow{\Psi}243\,\mathcal W_6
 \lhook\joinrel\longrightarrow\mathbb F_3^6,
 \qquad|\mathbb F_3^6|=729.
 \label{p12-83-c9-s1:eq:cadena-tipica-pi-autonoma}
\end{equation}
Los cardinales cuentan, respectivamente, condiciones iniciales
enriquecidas, emisiones distintas, palabras ternarias visibles y elementos
de la fibra ambiente completa.

Para \(U_6=(u_1,\ldots,u_6)\), la proyección es
\[
 \Psi(U_6)=((-u_1)\bmod3,\ldots,(-u_6)\bmod3).
\]
El censo de las \(43\) órbitas de \(D_3\) demuestra que la órbita de
clausura es la única cuyas seis orientaciones poseen a la vez:
\[
 \text{cinco elevaciones},\qquad
 \text{masa }1008,\qquad
 1008=432+4\cdot144.
\]
El calibre
\[
 \operatorname{diag}_c=6,
 \qquad
 \operatorname{card}=ES
\]
selecciona un único representante:
\begin{equation}
 \boxed{w_6^\pi=010211.}
 \label{p12-83-c9-s1:eq:semilla-pi-autonoma}
\end{equation}
La selección usa catálogo, multiplicidades, acción orbital y orientación;
no consulta un prefijo de \(\pi\).

\subsection{Elevaciones \(L_0,L_1\) y 1.ª frontera coinductiva}

Sobre vectores fila de \(\mathbb F_3^6\), sean
\[
L_0=
\begin{pmatrix}
2&2&2&1&2&1\\
2&1&2&2&1&1\\
1&0&0&1&0&2\\
0&0&1&1&1&0\\
2&2&2&0&2&1\\
2&1&2&1&0&0
\end{pmatrix},
\quad
L_1=
\begin{pmatrix}
2&1&2&0&2&2\\
1&2&1&1&2&0\\
1&2&1&1&1&2\\
0&1&0&0&2&1\\
2&0&2&1&0&1\\
0&2&2&2&1&1
\end{pmatrix}.
\]
Si \(b_1=w_6^\pi\), la recurrencia de bloques es
\[
 b_{k+1}=b_kL_{\varepsilon_k},
 \qquad
 (\varepsilon_1,\varepsilon_2,\varepsilon_3,\varepsilon_4)=(0,0,1,1),
\]
y la palabra acumulada se forma por concatenación:
\[
 w_{6(k+1)}=w_{6k}\mathbin{\|}b_{k+1}.
\]
No se multiplica una palabra de longitud \(6k\) por una matriz \(6\times6\).

Para el canal de clausura:
\[
\begin{array}{c|c}
\text{profundidad}&\text{bloque nuevo}\\ \hline
6&010211\\
12&012222\\
18&010211\\
24&002111\\
30&110221
\end{array}
\]
y
\begin{equation}
 \boxed{
 w_{30}^{\pi}
 =010211|012222|010211|002111|110221.}
 \label{p12-83-c9-s1:eq:w30-pi-autonomo}
\end{equation}
La frontera siguiente es
\[
 u_5^\pi=222220.
\]
Junto con los otros canales,
\[
 R_{36}=
 \begin{pmatrix}
 222220\\021222\\102011
 \end{pmatrix}.
\]
El símbolo \(R_{36}\) registra el sexto bloque de tres historias y abre una
frontera; no es un terminal.

La rama certificada de veinte bloques comienza
\begin{align*}
\pi:\;&010211|012222|010211|002111|110221|222220|111201|\\
&212121|200121|100100|101222|022212|012012|111210|\\
&121011|200220|120210|000101|022010|020111.
\end{align*}
Por ello, ni \(w_{30}\) ni \(R_{36}\) terminan la genealogía.

\subsection{Descomposición orientada de la microfibra de \texorpdfstring{\(\pi\)}{pi}: región central, \(4\) regiones laterales y multiplicidad \texorpdfstring{\(432+4\cdot144=1008\)}{432+4x144=1008}}

La microfibra
\[
 \mathscr R_\pi^{\rm micro}
 :=\{U\in\mathcal U_6:\Psi(U)=010211\}
\]
contiene exactamente cinco regiones:
\[
\begin{array}{c|c|c|c}
U_6&E_{\rm sig}&\operatorname{mult}&\text{papel}\\ \hline
501|614|498|169|272|272&555555&432&\perp\\
810|923|870|169|272|272&339555&144&++\\
810|983|810|169|272|272&393555&144&++\\
870|923|810|169|272|272&933555&144&++\\
870|923|810|418|674|521&933717&144&--
\end{array}
\]
En consecuencia,
\begin{equation}
 \boxed{
 5=1_\perp+3_{++}+1_{--},
 \qquad
 1008=432+4\cdot144.}
 \label{p12-83-c9-s1:eq:descomposiciones-cinco-regiones-pi}
\end{equation}
La primera identidad clasifica funciones de orientación. La segunda cuenta
multiplicidades de procedencia. No cuentan cinco copias indistinguibles.

La región transversal es
\[
 501|614|498|169|272|272,
 \qquad E_{\rm sig}=555555,
 \qquad\operatorname{mult}=432.
\]
El retorno mutado es
\[
 870|923|810|418|674|521,
 \qquad E_{\rm sig}=933717,
 \qquad\operatorname{mult}=144.
\]
Intercambiarlos conservaría el censo \(3/1/1\), pero falsaría la estructura
regional completa.

\subsection{Descenso de un bloque nonádico completo al multigrafo de transferencia de fase}

Para \(U_6=(U_1,\ldots,U_6)\), sea
\[
 A(U_6)=(U_1,U_2,U_3),
 \qquad
 B(U_6)=(U_4,U_5,U_6).
\]
Cada emisión define una arista entre \(A(U_6)\) y \(B(U_6)\). El multigrafo
bruto tiene \(96\) vértices, distribuidos en componentes de tamaños
\[
 28,28,28,4,4,4.
\]
Se cocienta la fase interna mediante
\[
 (a,b,c)\sim(b,c,a)\sim(c,a,b).
\]
El representante es la rotación lexicográficamente mínima. El grafo
cocientado posee \(32\) vértices y dos componentes de tamaños \(28\) y \(4\).

Los anclajes son
\[
\begin{aligned}
U_{\rm APP}&=903|850|850|252|109|252,\\
U_e&=850|963|890|149|252|212,\\
U_\varphi&=830|943|830|109|252|252.
\end{aligned}
\]
Numeramos las regiones
\[
\begin{aligned}
U_\pi^{(1)}&=810|923|870|169|272|272,\\
U_\pi^{(2)}&=810|983|810|169|272|272,\\
U_\pi^{(3)}&=870|923|810|169|272|272,\\
U_\pi^{(4)}&=870|923|810|418|674|521,\\
U_\pi^{(5)}&=501|614|498|169|272|272.
\end{aligned}
\]

\subsection{Los \(5\) ciclos orientados de clausura}

\subsubsection{Cierre de \texorpdfstring{\(U_\pi^{(1)}\)}{U-pi 1}}
\begin{align*}
&252|109|252|903|850|850\to
903|850|850|252|109|252\to\\
&252|109|252|923|870|810\to
810|923|870|169|272|272\to\\
&272|169|272|963|890|850\to
850|963|890|149|252|212\to\\
&212|149|252|943|830|830\to
830|943|830|109|252|252\to
252|109|252|903|850|850.
\end{align*}

\subsubsection{Cierre de \texorpdfstring{\(U_\pi^{(2)}\)}{U-pi 2}}
\begin{align*}
&903|850|850|252|109|252\to
272|169|272|903|850|850\to\\
&272|169|272|983|810|810\to
810|983|810|169|272|272\to\\
&272|169|272|963|890|850\to
850|963|890|149|252|212\to\\
&212|149|252|943|830|830\to
830|943|830|109|252|252\to
903|850|850|252|109|252.
\end{align*}

\subsubsection{Cierre de \texorpdfstring{\(U_\pi^{(3)}\)}{U-pi 3}}
\begin{align*}
&903|850|850|252|109|252\to
292|189|232|903|850|850\to\\
&292|189|232|923|810|870\to
870|923|810|169|272|272\to\\
&272|169|272|963|890|850\to
850|963|890|149|252|212\to\\
&212|149|252|943|830|830\to
830|943|830|109|252|252\to
903|850|850|252|109|252.
\end{align*}

\subsubsection{Cierre de \texorpdfstring{\(U_\pi^{(4)}\)}{U-pi 4}}
\begin{align*}
&903|850|850|252|109|252\to
272|169|272|903|850|850\to\\
&272|169|272|963|890|850\to
850|963|890|149|252|212\to\\
&212|149|252|923|810|870\to
870|923|810|418|674|521\to\\
&943|830|830|521|418|674\to
830|943|830|109|252|252\to
903|850|850|252|109|252.
\end{align*}

\subsubsection{Cierre de \texorpdfstring{\(U_\pi^{(5)}\)}{U-pi 5}}
\begin{align*}
&252|109|252|903|850|850\to
903|850|850|252|109|252\to\\
&614|498|501|252|109|252\to
501|614|498|169|272|272\to\\
&272|169|272|963|890|850\to
850|963|890|149|252|212\to\\
&212|149|252|943|830|830\to
830|943|830|109|252|252\to
252|109|252|903|850|850.
\end{align*}

\begin{theorem}[Minimalidad de los cinco cierres]
\label{p12-83-c9-s1:thm:minimalidad-cierres-pi-autonoma}
Para cada \(i\in\{1,\ldots,5\}\), la longitud mínima de un paseo cerrado
que recorre simultáneamente
\[
 U_\pi^{(i)},\quad U_e,\quad U_\varphi,\quad U_{\rm APP}
\]
en el multigrafo cocientado es ocho.
\end{theorem}

\begin{proof}
Fijado \(i\), sea
\[
 E_*=\{U_\pi^{(i)},U_e,U_\varphi,U_{\rm APP}\}.
\]
Se considera el espacio finito
\[
 \mathscr S_i=\{(v,M):v\in V(G),\ M\subseteq E_*\}.
\]
El primer componente registra el vértice y el segundo las aristas objetivo
ya recorridas. Si una arista \(e\) lleva \(v\) a \(v'\), la transición es
\[
 (v,M)\longmapsto
 \bigl(v',M\cup(E_*\cap\{e\})\bigr).
\]
La búsqueda en anchura enumera longitudes en orden creciente. El primer
retorno a \((v,E_*)\) aparece en longitud ocho para las cinco regiones.
Los paseos impresos realizan la cota; la optimalidad de la búsqueda excluye
una longitud menor.
\end{proof}

La función de clausura queda codificada por
\[
 \mathfrak C_\pi=
 \bigl(\mathscr R_\pi^{\rm micro},
 w_\pi,\rho_\pi,m_\pi,\ell_{\rm cl}\bigr),
\]
\[
 \rho_\pi=(\perp,++ ,++ ,++ ,--),
 \qquad
 m_\pi=(432,144,144,144,144),
 \qquad
 \ell_{\rm cl}=8.
\]

\subsection{Operador de incidencia y media vuelta orientada}

La sombra de incidencia es
\[
 A_W=
 \begin{pmatrix}
 0&1&1&1&1&1\\
 1&0&1&2&2&1\\
 1&1&0&1&2&2\\
 1&2&1&0&1&2\\
 1&2&2&1&0&1\\
 1&1&2&2&1&0
 \end{pmatrix}.
\]
La multiplicación en \(\mathbb F_3\) da
\begin{equation}
 \boxed{A_W^2=2I_6=-I_6.}
 \label{p12-83-c9-s1:eq:aw-cuadrado-menos-identidad-pi}
\end{equation}
Los productos entre filas y columnas distintas se anulan módulo tres,
mientras cada producto diagonal vale dos. Una aplicación realiza un cuarto
de giro ternario; dos aplicaciones realizan la inversión orientada.

\subsection{La pendiente 1/5 forzada por las \(5\) regiones}

Sea
\[
 \lambda=(\lambda_1,\ldots,\lambda_5)^{\mathsf T}
 \in\mathbb Q_{\geq0}^5,
 \qquad
 \mathbf1^{\mathsf T}\lambda=1.
\]
Olvidar la etiqueta regional aplica el simetrizador
\[
 P_5=\frac15\mathbf1\mathbf1^{\mathsf T}.
\]
Para todo vector normalizado,
\[
 P_5\lambda=\frac15\mathbf1.
\]
La invariancia \(P_5\lambda=\lambda\) tiene, por tanto, una única solución:
\begin{equation}
 \boxed{\lambda=\frac15\mathbf1,\qquad t_1=\frac15.}
 \label{p12-83-c9-s1:eq:pendiente-primitiva-pi}
\end{equation}
El entero cinco procede de la microfibra regional; no se reconoce en una
fórmula de \(\pi\).

\subsection{Composición racional y compensador \texorpdfstring{\(239\)}{239}}

La ley de composición de pendientes es
\[
 t\star s:=\frac{t+s}{1-ts}.
\]
La primera duplicación da
\[
 t_2=t_1\star t_1
 =\frac{2/5}{1-1/25}
 =\frac5{12}.
\]
La segunda da
\[
 t_4=t_2\star t_2
 =\frac{10/12}{1-25/144}
 =\frac{120}{119}.
\]
El retorno a la diagonal unitaria exige \(q>0\) tal que
\[
 \frac{120/119-1/q}{1+(120/119)/q}=1.
\]
Multiplicar por \(119q\) produce
\[
 120q-119=119q+120,
\]
y fuerza
\begin{equation}
 \boxed{q=239.}
 \label{p12-83-c9-s1:eq:compensador-239-pi}
\end{equation}

\subsection{Identidad gaussiana y 1.ª media vuelta}

Definimos
\begin{equation}
 \Theta_\pi
 :=16\arctan\frac15-4\arctan\frac1{239}.
 \label{p12-83-c9-s1:eq:theta-pi-autonoma}
\end{equation}
Los cuadrados sucesivos son
\[
\begin{array}{c|r@{}l}
\text{potencia}&\multicolumn{2}{c}{\text{entero gaussiano}}\\ \hline
(5+i)^2&24&+10i\\
(5+i)^4&476&+480i\\
(5+i)^8&-3824&+456960i\\
(5+i)^{16}&-208797818624&-3494830080i\\
(239-i)^2&57120&-478i\\
(239-i)^4&3262465916&-54606720i
\end{array}
\]
y la multiplicación final da
\begin{equation}
 \boxed{
 (5+i)^{16}(239-i)^4
 =-681386607803576157184.}
 \label{p12-83-c9-s1:eq:identidad-gaussiana-pi-autonoma}
\end{equation}
La parte imaginaria es cero y la real es negativa. Por tanto,
\[
 \Theta_\pi\equiv\pi\pmod{2\pi}.
\]
Para \(0<x<1\),
\[
 x-\frac{x^3}{3}<\arctan x<x.
\]
De aquí
\[
 \Theta_\pi>
 16\left(\frac15-\frac1{3\cdot5^3}\right)-\frac4{239}>3
\]
y
\[
 \Theta_\pi<\frac{16}{5}<4.
\]
La única orientación negativa en ese intervalo es la primera media vuelta
positiva:
\begin{equation}
 \boxed{\Theta_\pi=\pi.}
 \label{p12-83-c9-s1:eq:reconocimiento-pi-exacto}
\end{equation}

El orden del reconocimiento arquimediano posterior es
\[
 \text{cinco regiones}
 \to\frac15
 \to\frac5{12}
 \to\frac{120}{119}
 \to239
 \to\Theta_\pi
 \to\pi.
\]

\subsection{Horquillas racionales de precisión arbitraria}

Para \(0<x<1\), sea
\[
 S_n(x)=\sum_{r=0}^{n}(-1)^r\frac{x^{2r+1}}{2r+1}.
\]
El criterio de Leibniz da
\[
 S_{2m+1}(x)<\arctan x<S_{2m}(x).
\]
Definimos
\[
 a_m^-=S_{2m+1}(1/5),\quad a_m^+=S_{2m}(1/5),
\]
\[
 b_m^-=S_{2m+1}(1/239),\quad b_m^+=S_{2m}(1/239),
\]
y
\begin{equation}
 \ell_{\pi,m}=16a_m^- -4b_m^+,
 \qquad
 h_{\pi,m}=16a_m^+ -4b_m^-.
 \label{p12-83-c9-s1:eq:horquilla-pi-autonoma}
\end{equation}
Entonces
\[
 \ell_{\pi,m}<\pi<h_{\pi,m}
\]
y
\begin{equation}
 h_{\pi,m}-\ell_{\pi,m}
 =
 \frac{16}{(4m+3)5^{4m+3}}
 +\frac{4}{(4m+3)239^{4m+3}}
 \longrightarrow0.
 \label{p12-83-c9-s1:eq:anchura-horquilla-pi}
\end{equation}
Las cotas inferiores crecen y las superiores decrecen.

Como \(3<\pi<4\), el lector ternario actúa sobre
\[
 r_\pi:=\pi-3\in(0,1),
\]
con horquilla
\[
 [\ell_{\pi,m}-3,h_{\pi,m}-3].
\]

\subsection{Certificación arquimediana posterior de la fibra ya generada}

Supóngase que la relación interna \(\operatorname{Ext}\) y la supervivencia
coinductiva han producido ya, sin emplear Machin ni una expansión tabulada,
la extensión compatible \(N_{K+1}\) del prefijo \(N_K\), de longitud
\(L_K=6K\). Los \(729\) subcilindros de su fibra ambiente son
\[
 I_{K,u}=
 \left[
 \frac{729N_K+u}{3^{L_K+6}},
 \frac{729N_K+u+1}{3^{L_K+6}}
 \right),
 \qquad0\leq u<729.
\]
El orden \(m=m(K)\) se aumenta después hasta que la horquilla de Leibniz
localiza el subcilindro ya producido. Sean
\[
 \ell_{\pi,K}:=\ell_{\pi,m(K)}-3,
 \qquad
 h_{\pi,K}:=h_{\pi,m(K)}-3.
\]
Los extremos enteros compatibles son
\[
 j_{\min}=
 \max\!\left(
 729N_K,\left\lfloor
 \ell_{\pi,K}3^{L_K+6}\right\rfloor
 \right),
\]
\[
 j_{\max}=
 \min\!\left(
 729N_K+728,\left\lceil
 h_{\pi,K}3^{L_K+6}\right\rceil-1
 \right).
\]
La condición de localización unitaria es
\begin{equation}
 \boxed{j_{\min}=j_{\max},\qquad N_{K+1}=j_{\min}.}
 \label{p12-83-c9-s1:eq:seleccion-unitaria-subcilindro-pi}
\end{equation}
La primera igualdad certifica la localización unitaria; la segunda verifica
que esa localización coincide con la extensión generada internamente. Ninguna
de las dos define la transición TPK ni selecciona entre los \(729\) estados.
La existencia y la unicidad proceden de \(\operatorname{Ext}\) y de la
supervivencia; Machin y las horquillas de Leibniz son reconocimientos
posteriores.

\subsection{Prolongación coinductiva de la sección de \(\pi\) mediante el sistema inverso de historias supervivientes}

Sea \(\mathcal H_m^\pi\) el conjunto de historias enriquecidas a profundidad
\(m\), y sea
\[
 p_{n,m}:\mathcal H_n^\pi\longrightarrow\mathcal H_m^\pi
\]
el truncamiento. Definimos
\[
 \operatorname{Surv}_m^{(N),\pi}
 :=p_{N,m}(\mathcal H_N^\pi),
\]
\[
 \operatorname{Surv}_m^\pi
 :=\bigcap_{N\geq m}\operatorname{Surv}_m^{(N),\pi}.
\]
Cuando la relación \(\operatorname{Ext}\) es no vacía, univaluada y
compatible en todos los niveles, los conjuntos son unitarios y satisfacen
\[
 p_{m+1,m}(\operatorname{Surv}_{m+1}^{\pi})
 =\operatorname{Surv}_{m}^{\pi}.
\]
En ese dominio de prolongabilidad existe, por tanto, una única sección
\[
 x_\pi\in
 X_\pi:=\varprojlim_m\operatorname{Surv}_m^\pi.
\]
Los cilindros están encajados y sus diámetros son \(3^{-6K}\), que tiende a
cero. Su intersección contiene un único punto, identificado por el carácter
angular con \(r_\pi=\pi-3\).

Para toda profundidad decimal \(M\), puede elegirse un nivel \(K(M)\) cuyo
cilindro queda dentro de una sola celda decimal de diámetro \(10^{-M}\).
Los prefijos profundos truncan exactamente a los anteriores.

\subsection{Caracterización del canal coinductivo de \(\pi\): selección de \(w_6^\pi\), \(5\) regiones, inversión orientada y unicidad de la sección límite}

\begin{theorem}[Caracterización del canal de clausura]
\label{p12-83-c9-s1:thm:caracterizacion-final-pi-autonoma}
Partiendo del catálogo, la acción \(D_3\), el calibre orientado, los
operadores \(L_0,L_1,A_W\), la ley de pendientes y la filtración racional:
\begin{enumerate}
 \item se selecciona \(w_6^\pi=010211\);
 \item se construyen cinco regiones con las identidades de
       \eqref{p12-83-c9-s1:eq:descomposiciones-cinco-regiones-pi};
 \item cada región pertenece a un cierre mínimo de longitud ocho;
 \item \(A_W^2=-I_6\) fija la inversión orientada;
 \item la simetrización regional fuerza \(1/5\);
 \item la composición fuerza \(120/119\);
 \item el retorno unitario fuerza \(239\);
 \item la identidad gaussiana reconoce la primera media vuelta;
 \item las horquillas y la fibra de \(729\) elementos producen una sección
       compatible a toda profundidad finita.
\end{enumerate}
Su evaluación es
\[
 \boxed{\pi=16\arctan\frac15-4\arctan\frac1{239}.}
\]
\end{theorem}

La cadena completa es
\[
 \mathrm{APP}\to\mathrm{TRIT}\to\mathrm{TPK}
 \to104\,976\to468\to243\to w_6^\pi
 \to w_{30}^\pi\to R_{36}\to\Gamma_9
 \to X_\pi\to\pi.
\]

\begin{criterion}[Falsación del canal de clausura]
La derivación queda refutada si:
\begin{enumerate}
 \item la órbita de clausura deja de ser única;
 \item desaparece una de las cinco regiones o cambia su identidad completa;
 \item uno de los cinco cierres tiene longitud menor que ocho;
 \item \(A_W^2\ne2I_6\);
 \item el simetrizador admite otra distribución normalizada;
 \item las composiciones no producen \(5/12\), \(120/119\) y \(239\);
 \item falla la identidad gaussiana;
 \item una horquilla deja de contener el carácter angular;
 \item un nivel no recorre los \(729\) elementos de la fibra ambiente o no
       localiza exactamente una extensión;
 \item una expansión objetivo interviene antes del reconocimiento.
\end{enumerate}
\end{criterion}
\fi

La construcción completa del carácter de clausura, las cinco regiones y la
biografía estructural de \(\pi\) quedó establecida en el
\cref{chap:83-08}. Partiendo de ese mismo carácter ya construido, el presente
capítulo prolonga ahora el refinamiento hacia \(e\) y \(\varphi\), sin
reiniciar la genealogía ni repetir su primera etapa.


% Unidad U013. Biografía estructural completa del carácter exponencial.

\section{Propagación, retorno con memoria y construcción del carácter exponencial}
\label{p12-83-c9-s2:chap:biografia-completa-e}

\subsection{Selector orbital finito: existencia y unicidad de la emisión compatible}

Sea \(\mathcal W=\mathbb F_3^6\). Para una emisión
\(U=(u_1,\ldots,u_6)\), la palabra visible es
\[
 \Psi(U)=((-u_1)\bmod3,\ldots,(-u_6)\bmod3).
\]
Sobre \(\mathcal W\), la acción diedral está dada por
\[
\begin{aligned}
 r(u_1,u_2,u_3,u_4,u_5,u_6)
 &=(u_3,u_1,u_2,u_5,u_6,u_4),\\
 s(u_1,u_2,u_3,u_4,u_5,u_6)
 &=(u_4,u_5,u_6,u_1,u_2,u_3).
\end{aligned}
\]
Para \(w\in\mathcal W\), definimos
\[
 n(w)=\#\Psi^{-1}(w),
 \qquad
 M(w)=\sum_{U\in\Psi^{-1}(w)}\operatorname{mult}(U),
\]
y
\[
 \nu(w)=
 \bigl(
 \#\{j:w_j=0\},
 \#\{j:w_j=1\},
 \#\{j:w_j=2\}
 \bigr).
\]

\begin{theorem}[Selección estructural del canal exponencial]
\label{p12-83-c9-s2:thm:seleccion-estructural-e}
Entre las \(43\) órbitas del catálogo visible, existe una única órbita cuyas
fibras satisfacen simultáneamente:
\[
 n(w)=1,\qquad M(w)=144,\qquad
 \mathfrak D(w)=\{0,3,6\},\qquad
 \nu(w)=(2,3,1).
\]
El calibre
\[
 \operatorname{diag}_c=6,
 \qquad \operatorname{card}=ES
\]
selecciona en ella
\[
 \boxed{w_6^e=201101.}
\]
\end{theorem}

\begin{proof}
Las cuatro condiciones se evalúan sobre las \(243\) palabras del catálogo,
no sobre una muestra. Una única órbita satisface el predicado conjunto. En
esa órbita, una sola emisión posee a la vez diagonal \(6\) y orientación
\(ES\). La enumeración usa únicamente campos del catálogo APP--TRIT y no
consulta una expansión de \(e\).
\end{proof}

\subsection{Elevación visible y cadena de bloques}

La recurrencia \(L_0,L_0,L_1,L_1\) produce
\[
 w_{30}^{e}
 =201101|121221|102011|012222|102011,
\]
y la frontera conjunta aporta
\[
 u_5^e=021222.
\]
El registro de treinta y seis trits queda
\[
 w_{36}^{e}
 =201101|121221|102011|012222|102011|021222.
\]

\subsection{Firma decimal y cuadrado posicional}

La lectura inicial de doce cifras es
\[
 \mathbf d_e^{(12)}
 =718281828459
 =7\,\underbrace{1828\,1828}_{\text{cuadrado local}}\,459.
\]
La igualdad central es la concatenación
\[
 1828\mathbin{\|}1828,
\]
no un producto aritmético ni el comienzo supuesto de un período.

\begin{definition}[Cuadrado posicional]
Una palabra \(d_1\cdots d_{12}\) contiene un cuadrado de longitud \(\ell\)
en la posición \(p\) cuando
\[
 d_p\cdots d_{p+\ell-1}
 =d_{p+\ell}\cdots d_{p+2\ell-1}.
\]
La posición y los contextos laterales pertenecen a la firma.
\end{definition}

\begin{table}[htbp]
\centering
\caption{Censo exhaustivo de cuadrados en las lecturas de doce cifras.}
\label{p12-83-c9-s2:tab:censo-cuadrados-e}
\begin{tabular}{c*{6}{r}}
\toprule
Longitud \(\ell\)&1&2&3&4&5&6\\
\midrule
Ocurrencias&240&21&3&2&0&0\\
Palabras que las contienen&153&19&3&2&0&0\\
\bottomrule
\end{tabular}
\end{table}

La otra palabra con un cuadrado de longitud cuatro es
\[
 575157518472
 =\underbrace{5751\,5751}_{\text{posición inicial }1}\,8472.
\]

\begin{proposition}[Unicidad posicional]
\label{p12-83-c9-s2:prop:cuadrado-posicional-e}
La lectura del canal \(e\) es la única que posee un cuadrado de longitud
cuatro después de un prefijo de longitud uno y con contextos \(7\) y \(459\).
\end{proposition}

\begin{proof}
El censo deja sólo dos palabras con cuadrados de longitud cuatro. Una lo
inicia en la primera posición; la otra después del prefijo \(7\). La
posición y los contextos separan las dos posibilidades.
\end{proof}

\subsection{Retorno visible y ruptura de memoria}

En la elevación
\[
 201101|121221|102011|012222|102011
\]
el tercer y el quinto bloques coinciden:
\[
 B_3=B_5=102011.
\]
Sus preestados módulo \(27\) son, sin embargo,
\[
 p_3=(25,11,7,17,8,16),
 \qquad
 p_5=(7,8,4,23,26,7),
\]
y los cocientes posteriores:
\[
 q_3=(6,13,9,2,13,1),
 \qquad
 q_5=(15,0,3,19,23,25).
\]
Se verifican \(p_3\ne p_5\) y \(q_3\ne q_5\).

\begin{figure}[htbp]
\centering
\begin{tikzpicture}[
 node distance=11mm and 7mm,
 every node/.style={font=\small},
 block/.style={draw,rounded corners,minimum width=17mm,minimum height=8mm},
 state/.style={draw,dashed,rounded corners,align=center,inner sep=3pt},
 >=stealth
]
\node[block] (b1) {\(201101\)};
\node[block,right=of b1] (b2) {\(121221\)};
\node[block,right=of b2] (b3) {\(102011\)};
\node[block,right=of b3] (b4) {\(012222\)};
\node[block,right=of b4] (b5) {\(102011\)};
\draw[->] (b1)--(b2);
\draw[->] (b2)--(b3);
\draw[->] (b3)--(b4);
\draw[->] (b4)--(b5);
\node[state,below=of b3] (p3)
 {\(p_3=(25,11,7,17,8,16)\)\\
  \(q_3=(6,13,9,2,13,1)\)};
\node[state,below=of b5] (p5)
 {\(p_5=(7,8,4,23,26,7)\)\\
  \(q_5=(15,0,3,19,23,25)\)};
\draw[->] (p3)--(b3);
\draw[->] (p5)--(b5);
\draw[<->,thick] (b3) to[bend left=28]
 node[above]{misma palabra visible} (b5);
\end{tikzpicture}
\caption{Retorno visible de \(102011\) con renovación del preestado y del
cociente.}
\label{p12-83-c9-s2:fig:ruptura-memoria-e}
\end{figure}

\begin{theorem}[Retorno visible sin periodicidad del estado]
\label{p12-83-c9-s2:thm:retorno-visible-no-periodico-e}
\(B_3=B_5\) es un retorno de la proyección visible, no una órbita periódica
del estado enriquecido.
\end{theorem}

\begin{proof}
El retorno del estado completo exigiría igualdad de todas sus coordenadas.
Las desigualdades de preestado y cociente lo excluyen. Tampoco
\(1828|1828\) inicia un período decimal: después del segundo bloque aparece
\(4\), mientras una continuación periódica exigiría \(1\).
\end{proof}

\subsection{Certificación factorial posterior de la sección exponencial coinductiva}

Sea \(\mathcal A\) un álgebra conmutativa unitaria sobre \(\mathbb Q\), y
sea \(P_t\in\mathcal A[[t]]\) el propagador formal. La concatenación
temporal impone
\[
 P_{s+t}=P_sP_t,
 \qquad
 P_0=1.
\]
Escribimos
\[
 P_t=\sum_{n=0}^{\infty}a_nt^n.
\]
La unidad fija \(a_0=1\), y el transporte elemental normalizado fija
\(a_1=1\).

\begin{theorem}[Recurrencia factorial]
\label{p12-83-c9-s2:thm:recurrencia-factorial-e}
\[
 (n+1)a_{n+1}=a_n
 \qquad(n\geq0),
\]
y, por tanto,
\[
 a_n=\frac1{n!},
 \qquad
 P_t=\sum_{n=0}^{\infty}\frac{t^n}{n!}.
\]
En \(t=1\),
\[
 \boxed{P_1=e.}
\]
\end{theorem}

\begin{proof}
El coeficiente de \(s^nt\) en \(P_{s+t}\) es
\[
 \binom{n+1}{n}a_{n+1}=(n+1)a_{n+1}.
\]
En \(P_sP_t\) es \(a_na_1=a_n\). La igualdad de series produce la
recurrencia. Una inducción desde \(a_0=1\) da \(a_n=1/n!\).
\end{proof}

\begin{corollary}[Ecuación diferencial]
\[
 \frac{d}{dt}P_t=P_t,
 \qquad P_0=1.
\]
\end{corollary}

\begin{proof}
\[
 P_t'=\sum_{n\geq0}(n+1)a_{n+1}t^n
 =\sum_{n\geq0}a_nt^n=P_t.
\]
\end{proof}

\subsection{Encaje cofinal de horquillas racionales y convergencia al valor publicado}

Sea
\[
 S_N=\sum_{n=0}^{N}\frac1{n!}.
\]

\begin{theorem}[Cota factorial explícita]
\label{p12-83-c9-s2:thm:cota-factorial-e}
Para \(N\geq1\),
\begin{equation}
 \boxed{
 S_N<e<
 S_N+\frac{N+2}{(N+1)(N+1)!}.}
 \label{p12-83-c9-s2:eq:horquilla-factorial-e}
\end{equation}
\end{theorem}

\begin{proof}
El resto es
\[
 R_N=\sum_{k=0}^{\infty}\frac1{(N+1+k)!}>0.
\]
Para \(k\geq0\),
\[
 (N+1+k)!
 \geq(N+1)!(N+2)^k.
\]
Por tanto,
\[
 R_N
 \leq\frac1{(N+1)!}
 \sum_{k=0}^{\infty}\frac1{(N+2)^k}
 =\frac{N+2}{(N+1)(N+1)!}.
\]
La desigualdad es estricta desde \(k=2\).
\end{proof}

Definimos
\[
 J_N^{(e)}=
 \left[
 S_N,\,
 S_N+\frac{N+2}{(N+1)(N+1)!}
 \right]
\]
y, para la parte fraccionaria,
\[
 \widetilde J_N^{(e)}=J_N^{(e)}-2.
\]
La anchura converge a cero. Para cada \(M\), existe \(N\) tal que
\[
 \operatorname{diam}J_N^{(e)}<10^{-M}.
\]
Si los extremos pertenecen a una misma celda decimal, las primeras \(M\)
cifras quedan certificadas mediante aritmética racional.

\subsection{Prolongación coinductiva de la sección exponencial mediante subcilindros ternarios del Topological Phase Kernel}

En el nivel \(K\), TPK y la relación \(\operatorname{Ext}\) han generado ya
un único subcilindro compatible. Se toma el menor \(N=N_e(K)\) para el cual
la horquilla factorial queda contenida en ese subcilindro. \(N_e(K)\) es un
orden de certificación analítica, no un índice generador ni una regla de
selección de la fibra.

Los cilindros TPK son compatibles bajo truncamiento con independencia del
semigrupo. El teorema semigrupal identifica después su evaluación
arquimediana con \(e-2\); la parte entera recupera \(e\). Si fallan la órbita,
\(\operatorname{Ext}\) o el truncamiento, falla la generación; si fallan la
recurrencia o la cota factorial, falla este certificado analítico, no la
genealogía APP--TRIT--TPK.

\subsection{Certificado adversarial de la generación de \(e\): unicidad orbital, recurrencia factorial, horquillas racionales y prolongación cilíndrica}

\begin{criterion}[Falsación del canal exponencial]
La construcción queda refutada si:
\begin{enumerate}
 \item otra órbita satisface el predicado finito completo;
 \item la firma \(7[1828\,1828]459\) no es posicionalmente única;
 \item coinciden los preestados o cocientes de los dos bloques \(102011\);
 \item el semigrupo no fuerza \((n+1)a_{n+1}=a_n\);
 \item la cota \eqref{p12-83-c9-s2:eq:horquilla-factorial-e} falla;
 \item un nivel no localiza de forma unitaria una extensión compatible;
 \item una expansión tabulada interviene antes del reconocimiento.
\end{enumerate}
\end{criterion}


% Unidad U014. Biografía estructural completa del carácter de autoescala.

\section{Incidencia modal, rayo de Perron y construcción del carácter de autoescala}
\label{p12-83-c9-s3:chap:biografia-completa-phi}

\subsection{Selección finita de la semilla de autoescala}

Sea \(\mathcal U_6\subset\{0,\ldots,999\}^6\) el catálogo de las
\(468\) emisiones distintas construido en el
\cref{chap:semillas-emisor-54}, y sea
\[
 \Psi:\mathcal U_6\longrightarrow\mathbb F_3^6,
 \qquad
 \Psi(U)=((-U_j)\bmod3)_{j=1}^{6}.
\]
Para \(w\in\Psi(\mathcal U_6)\), recordamos los invariantes
\[
 n(w)=\#\Psi^{-1}(w),
 \qquad
 M(w)=\sum_{U\in\Psi^{-1}(w)}\operatorname{mult}(U),
\]
y el vector de ocupación
\[
 \nu(w)=
 \bigl(\#\{j:w_j=0\},\#\{j:w_j=1\},\#\{j:w_j=2\}\bigr).
\]

\begin{theorem}[Selección orbital de la autoescala]
\label{p12-83-c9-s3:thm:seleccion-orbital-phi}
La acción diedral sobre \(\Psi(\mathcal U_6)\) contiene una única órbita
que satisface el predicado de autoescala
\[
 n(w)=1,
 \qquad
 M(w)=144,
 \qquad
 \nu(w)=(2,2,2).
\]
El calibre interno
\[
 \operatorname{diag}_c=6,
 \qquad
 \operatorname{card}=ES
\]
selecciona en esa órbita la palabra
\[
 \boxed{w_6^{\varphi}=121200.}
\]
Su única emisión es
\[
 \boxed{U_{\varphi}=(830,943,830,109,252,252),}
\]
de multiplicidad \(144\), firma multiplicativa \(555933\), orientación
aditiva \(ES\), soporte multiplicativo de cardinal seis y tipo residual
\(A\).
\end{theorem}

\begin{proof}
La proyección directa da
\[
 U_{\varphi}\bmod3=(2,1,2,1,0,0),
 \qquad
 -U_{\varphi}\bmod3=(1,2,1,2,0,0).
\]
La enumeración se efectúa sobre las \(243\) palabras visibles y sus
\(43\) órbitas, no sobre una selección de ejemplos. El predicado conjunto
\(n=1\), \(M=144\), \(\nu=(2,2,2)\) deja una órbita. Dentro de ella, la
diagonal y la orientación declaradas dejan una emisión. En esta etapa no
se ha introducido la ecuación \(x^2-x-1=0\), una tabla decimal ni el nombre
convencional de la constante.
\end{proof}

\subsection{Elevación por bloques y 1.ª frontera coinductiva}

Sea \(b_1=w_6^{\varphi}\), y sea
\(\varepsilon=(0,0,1,1)\). Con las matrices \(L_0,L_1\) construidas en el
\cref{hmtctx:p12:chap:orbitas-elevacion-r36}, definimos
\[
 b_{k+1}=b_kL_{\varepsilon_k}\pmod3
 \qquad(1\leq k\leq4),
\]
y concatenamos
\[
 w_{6(k+1)}^{\varphi}=w_{6k}^{\varphi}\mathbin{\|}b_{k+1}.
\]
El cálculo matricial produce, bloque por bloque,
\[
\begin{aligned}
 b_1&=121200,\\
 b_2&=112202,\\
 b_3&=121020,\\
 b_4&=010210,\\
 b_5&=010200,
\end{aligned}
\]
y, por tanto,
\begin{equation}
 \boxed{
 w_{30}^{\varphi}
 =121200|112202|121020|010210|010200.}
 \label{p12-83-c9-s3:eq:w30-phi}
\end{equation}
La neutralidad dual de la frontera añade
\[
 u_5^{\varphi}=102011,
\]
de modo que
\begin{equation}
 \boxed{
 w_{36}^{\varphi}
 =121200|112202|121020|010210|010200|102011.}
 \label{p12-83-c9-s3:eq:w36-phi}
\end{equation}

\begin{remark}[Dos usos distintos de una palabra visible]
La palabra \(102011\) es también visible en la biografía de \(e\). La
igualdad de una proyección no identifica los estados enriquecidos: el canal,
el prefijo, el cociente, el acarreo, la fase y la memoria pertenecen al tipo
del objeto. Esta separación es obligatoria antes de construir cualquier
carácter real.
\end{remark}

\subsection{Del calendario trítico al multigrafo de incidencia}

\subsubsection{Incidencia modal primitiva \(F_{\rm av}\) y dinámica de memoria de orden \(1\)}

El calendario trítico primitivo es
\[
 (+1,-1,0).
\]
Fijamos dos modos visibles, \(\Sigma\) y \(\Pi\), y la regla
\begin{equation}
 +1:m\longmapsto\Sigma,
 \qquad
 -1:m\longmapsto\Pi,
 \qquad
 0:m\longmapsto m.
 \label{p12-83-c9-s3:eq:regla-modal-primitiva-phi}
\end{equation}
Con la condición de cierre cíclico, la órbita modal producida por
\eqref{p12-83-c9-s3:eq:regla-modal-primitiva-phi} es
\[
 \boxed{(\Sigma,\Pi,\Pi).}
\]
Los tres pares consecutivos son
\[
 \Sigma\longrightarrow\Pi,
 \qquad
 \Pi\longrightarrow\Pi,
 \qquad
 \Pi\longrightarrow\Sigma.
\]

\begin{definition}[Matriz de incidencia cronológica]
Para \(i,j\in\{\Sigma,\Pi\}\), sea
\[
 N_3(i,j)=
 \#\{r\in\mathbb Z/3\mathbb Z:(m_r,m_{r+1})=(i,j)\}.
\]
El subíndice \(3\) cuenta instantes del calendario primitivo; no denota una
potencia matricial.
\end{definition}

\begin{theorem}[Descenso necesario de la incidencia modal]
\label{p12-83-c9-s3:thm:descenso-incidencia-modal-phi}
En el orden \((\Sigma,\Pi)\),
\begin{equation}
 \boxed{
 N_3=
 \begin{pmatrix}
 0&1\\
 1&1
 \end{pmatrix}.}
 \label{p12-83-c9-s3:eq:N3-phi}
\end{equation}
Por repetición exacta del calendario,
\begin{equation}
 \boxed{
 N_9=3N_3,
 \qquad
 N_{27}=9N_3,
 \qquad
 N_{108}=36N_3.}
 \label{p12-83-c9-s3:eq:incidencias-9-27-108-phi}
\end{equation}
Estas tres igualdades son igualdades de conteos de aristas; no afirman
\(N_9=N_3^3\), \(N_{27}=N_3^9\) ni \(N_{108}=N_3^{36}\).
\end{theorem}

\begin{proof}
No aparece \(\Sigma\to\Sigma\); cada uno de los otros tres pares aparece
una vez. Esto determina las cuatro entradas de \(N_3\). Los calendarios de
longitud \(9\), \(27\) y \(108\) contienen respectivamente \(3\), \(9\)
y \(36\) copias del bloque primitivo. Cada repetición multiplica todos los
conteos por el mismo entero, lo que prueba
\eqref{p12-83-c9-s3:eq:incidencias-9-27-108-phi}.
\end{proof}

\subsubsection{Publicación areal–volumétrica y razón estable \(\varphi^{-2}\) bajo la medida de Parry}

La etiqueta modal y la hoja geométrica son objetos distintos. La carta
de publicación es
\[
 \mathfrak g_{\rm av}(\Sigma)=\mathscr H_{\rm ar},
 \qquad
 \mathfrak g_{\rm av}(\Pi)=\mathscr H_{\rm vol}.
\]
Transportando \eqref{p12-83-c9-s3:eq:N3-phi} por \(\mathfrak g_{\rm av}\), obtenemos
el operador de incidencia
\begin{equation}
 \boxed{
 F_{\rm av}=
 \begin{pmatrix}
 0&1\\
 1&1
 \end{pmatrix}.}
 \label{p12-83-c9-s3:eq:Fav-derivada-phi}
\end{equation}
Así, las dos transiciones cruzadas, la única autorretención volumétrica y la
ausencia de autorretención areal se derivan del calendario; no son cuatro
entradas independientes.

\begin{remark}[Alcance del descenso]
Olvidar la fase elemental no produce por sí solo un cociente de Markov
fuerte del Topological Phase Kernel, porque el modo visible no determina
cuál de los tres operadores internos actuará en el siguiente paso. Lo que
sí desciende sin hipótesis adicional es el multigrafo de aristas de un
bloque completo. Su envolvente mínimo de memoria uno tiene matriz de
adyacencia \(F_{\rm av}\).
\end{remark}

\subsection{Rayo positivo invariante y unicidad de la autoescala}

\begin{theorem}[Rayo de Perron derivado]
\label{p12-83-c9-s3:thm:rayo-perron-phi}
El cono positivo de \(F_{\rm av}\) contiene un único rayo propio. Si lo
normalizamos como \(v=(1,x)^{\mathsf T}\), entonces
\[
 x^2-x-1=0,
 \qquad
 x>0,
\]
y, por tanto,
\begin{equation}
 \boxed{x=\frac{1+\sqrt5}{2}=\varphi.}
 \label{p12-83-c9-s3:eq:phi-perron}
\end{equation}
\end{theorem}

\begin{proof}
La ecuación \(F_{\rm av}v=\lambda v\) da, coordenada a coordenada,
\[
 x=\lambda,
 \qquad
 1+x=\lambda x.
\]
Al eliminar \(\lambda\) resulta \(x^2-x-1=0\). Sus raíces son
\((1\pm\sqrt5)/2\), y sólo la positiva pertenece al interior del cono. La
matriz \(F_{\rm av}\) es primitiva, pues \(F_{\rm av}^2\) tiene todas sus
entradas positivas; el teorema de Perron--Frobenius garantiza que ese rayo
positivo es simple y único. El nombre \(\varphi\) se asigna después de esta
construcción.
\end{proof}

\begin{corollary}[Ecuación escalar de autoescala]
\[
 \varphi=1+\frac1\varphi,
 \qquad
 \varphi^2=\varphi+1,
 \qquad
 \varphi^{-1}=\varphi-1.
\]
Estas identidades son consecuencias de \eqref{p12-83-c9-s3:eq:phi-perron}; no seleccionan
la semilla ni la matriz.
\end{corollary}

\subsection{Recurrencia de Fibonacci del rayo de Perron y encaje diofántico de las horquillas racionales}

Sea \(F_0=0\), \(F_1=1\) y
\[
 F_{n+1}=F_n+F_{n-1}\qquad(n\geq1).
\]

\begin{theorem}[Potencias exactas de la incidencia]
\label{p12-83-c9-s3:thm:potencias-fibonacci-phi}
Para todo \(n\geq1\),
\begin{equation}
 \boxed{
 F_{\rm av}^{\,n}=
 \begin{pmatrix}
 F_{n-1}&F_n\\
 F_n&F_{n+1}
 \end{pmatrix}.}
 \label{p12-83-c9-s3:eq:potencias-Fav-fibonacci}
\end{equation}
\end{theorem}

\begin{proof}
Para \(n=1\), la identidad es la definición de \(F_{\rm av}\). Si vale
para \(n\), la multiplicación a la derecha por \(F_{\rm av}\) da
\[
 \begin{pmatrix}
 F_n&F_{n-1}+F_n\\
 F_{n+1}&F_n+F_{n+1}
 \end{pmatrix}
 =
 \begin{pmatrix}
 F_n&F_{n+1}\\
 F_{n+1}&F_{n+2}
 \end{pmatrix},
\]
que es el caso \(n+1\).
\end{proof}

\begin{theorem}[Identidad de Cassini]
\label{p12-83-c9-s3:thm:cassini-phi}
Para todo \(n\geq1\),
\begin{equation}
 \boxed{F_{n+1}F_{n-1}-F_n^2=(-1)^n.}
 \label{p12-83-c9-s3:eq:cassini-phi}
\end{equation}
\end{theorem}

\begin{proof}
El determinante de \(F_{\rm av}\) es \(-1\). Tomando determinantes en
\eqref{p12-83-c9-s3:eq:potencias-Fav-fibonacci},
\[
 F_{n-1}F_{n+1}-F_n^2
 =\det(F_{\rm av}^{\,n})=(-1)^n.
\]
\end{proof}

Definimos \(r_n=F_{n+1}/F_n\) para \(n\geq1\). De Cassini se sigue
\[
 r_n-r_{n-1}
 =\frac{F_{n+1}F_{n-1}-F_n^2}{F_nF_{n-1}}
 =\frac{(-1)^n}{F_nF_{n-1}}.
\]
Por tanto, las aproximaciones alternan alrededor del rayo de Perron.

\begin{proposition}[Horquillas racionales de Perron]
\label{p12-83-c9-s3:prop:horquillas-perron-phi}
Para \(m\geq1\),
\begin{equation}
 \frac{F_{2m+2}}{F_{2m+1}}
 <\varphi<
 \frac{F_{2m+1}}{F_{2m}},
 \label{p12-83-c9-s3:eq:horquilla-fibonacci-phi}
\end{equation}
y la anchura es
\begin{equation}
 \boxed{
 \frac{F_{2m+1}}{F_{2m}}
 -\frac{F_{2m+2}}{F_{2m+1}}
 =\frac{1}{F_{2m}F_{2m+1}}.}
 \label{p12-83-c9-s3:eq:anchura-fibonacci-phi}
\end{equation}
\end{proposition}

\begin{proof}
La alternancia anterior sitúa los cocientes pares por debajo y los impares
por encima de \(\varphi\). La fórmula de la anchura es
\eqref{p12-83-c9-s3:eq:cassini-phi} con \(n=2m+1\), tras poner ambas fracciones sobre
denominador común. Como \(F_n\to\infty\), la anchura tiende a cero.
\end{proof}

\subsection{Incidencia de memoria de orden \(1\) y medida de Parry del envolvente simbólico}

Sea \(X_{\rm av}\) el menor subshift de tipo finito sobre
\(\{\mathscr H_{\rm ar},\mathscr H_{\rm vol}\}\) que admite exactamente
los tres pares observados en el calendario modal. Su matriz de adyacencia es
\(F_{\rm av}\).

\begin{theorem}[Transición y medida de Parry]
\label{p12-83-c9-s3:thm:parry-phi}
Sea \(r=(1,\varphi)^{\mathsf T}\). La matriz estocástica de Parry es
\begin{equation}
 \boxed{
 P_{ij}=
 \frac{(F_{\rm av})_{ij}r_j}{\varphi r_i}
 =
 \begin{pmatrix}
 0&1\\
 \varphi^{-2}&\varphi^{-1}
 \end{pmatrix}.}
 \label{p12-83-c9-s3:eq:parry-matriz-phi}
\end{equation}
La medida estacionaria de vértices es proporcional a
\((1,\varphi^2)\), luego
\begin{equation}
 \boxed{
 \frac{\mu_{\rm ar}}{\mu_{\rm vol}}=\varphi^{-2}.}
 \label{p12-83-c9-s3:eq:parry-razon-phi}
\end{equation}
\end{theorem}

\begin{proof}
La suma de la fila \(i\) de \(P\) es
\[
 \frac{(F_{\rm av}r)_i}{\varphi r_i}=1.
\]
Como \(F_{\rm av}\) es simétrica, los vectores propios izquierdo y derecho
coinciden. Los pesos estacionarios son proporcionales a sus productos
coordenada a coordenada, es decir, a \((1,\varphi^2)\). La razón entre los
dos pesos es \(\varphi^{-2}\).
\end{proof}

\begin{remark}[Tres medidas que no deben confundirse]
La frecuencia cronológica de la órbita primitiva es
\[
 \nu_{\rm cyc}(\mathscr H_{\rm ar})=\frac13,
 \qquad
 \nu_{\rm cyc}(\mathscr H_{\rm vol})=\frac23.
\]
La medida uniforme \(\tau_{468}\) vive sobre el catálogo de emisiones. La
medida de Parry vive sobre historias admisibles de longitud arbitraria. La
primera cuenta instantes, la segunda cuenta emisiones y la tercera pondera
cilindros dinámicos. Ninguna es una aproximación de las otras.
\end{remark}

\subsection{Carácter multiplicativo y conservación cilíndrica}

Para una ruta admisible \(\gamma:i\to j\) de longitud \(n\), definimos
\begin{equation}
 \chi_P(\gamma)=\varphi^n\frac{r_i}{r_j}.
 \label{p12-83-c9-s3:eq:caracter-parry-phi}
\end{equation}

\begin{proposition}[Multiplicatividad por concatenación]
\label{p12-83-c9-s3:prop:multiplicatividad-parry-phi}
Si \(\gamma_1:i\to j\) y \(\gamma_2:j\to k\), entonces
\[
 \chi_P(\gamma_2\circ\gamma_1)
 =\chi_P(\gamma_2)\chi_P(\gamma_1).
\]
Sobre una ruta cerrada de longitud \(n\),
\[
 \chi_P(\gamma)=\varphi^n.
\]
\end{proposition}

\begin{proof}
La longitud se suma y la coordenada intermedia telescopa:
\[
 \varphi^{n_1+n_2}\frac{r_i}{r_k}
 =\left(\varphi^{n_1}\frac{r_i}{r_j}\right)
  \left(\varphi^{n_2}\frac{r_j}{r_k}\right).
\]
Si \(i=k\), el cociente terminal es uno.
\end{proof}

Normalizamos el vector izquierdo como
\[
 \ell=\frac{r}{1+\varphi^2},
 \qquad
 \ell^{\mathsf T}r=1.
\]
La medida de un cilindro es
\[
 \mu[x_0\cdots x_n]
 =\ell_{x_0}r_{x_n}\varphi^{-n}.
\]
Por tanto,
\begin{equation}
 V_n(x)
 :=\frac{\mu[x_0]}{\mu[x_0\cdots x_n]}
 =\varphi^n\frac{r_{x_0}}{r_{x_n}}
 =\chi_P(x_0\cdots x_n).
 \label{p12-83-c9-s3:eq:volumen-cilindrico-phi}
\end{equation}
Para una dimensión tipada \(D>0\), la escala
\[
 a_n^{(D)}(x)=V_n(x)^{1/D}
\]
satisface la ley exacta
\begin{equation}
 \boxed{
 \bigl(a_n^{(D)}(x)\bigr)^D
 \mu[x_0\cdots x_n]=\mu[x_0].}
 \label{p12-83-c9-s3:eq:conservacion-cilindrica-phi}
\end{equation}

En dimensión tres y sobre retornos cerrados,
\[
 \frac{a_9}{a_0}=\varphi^3,
 \qquad
 \frac{a_{27}}{a_0}=\varphi^9,
 \qquad
 \frac{a_{108}}{a_0}=\varphi^{36}.
\]
Si \(A_k=a_{9k}\), entonces
\begin{equation}
 A_{k+1}-2A_k+A_{k-1}
 =A_{k-1}(\varphi^3-1)^2>0.
 \label{p12-83-c9-s3:eq:convexidad-retornos-phi}
\end{equation}
Ésta es convexidad en profundidad de retorno; no se interpreta como una
magnitud cinemática sin una carta adicional de duración.

\subsection{Memoria cuadrática y rama contraída}

Los autovalores de \(F_{\rm av}\) son
\[
 \varphi,
 \qquad
 -\varphi^{-1}.
\]
La rama contraída cambia de orientación en cada iteración. Para conservar
ese signo antes de elevar al cuadrado, pasamos a
\(\operatorname{Sym}^2(\mathbb R^2)\). En la base
\((x^2,2xy,y^2)\),
\begin{equation}
 \operatorname{Sym}^2(F_{\rm av})=
 \begin{pmatrix}
 0&0&1\\
 0&1&2\\
 1&1&1
 \end{pmatrix}.
 \label{p12-83-c9-s3:eq:sym2-Fav-phi}
\end{equation}

\begin{theorem}[Espectro de la memoria de segundo orden]
\label{p12-83-c9-s3:thm:espectro-sym2-phi}
\[
 \det\bigl(tI-\operatorname{Sym}^2(F_{\rm av})\bigr)
 =(t+1)(t^2-3t+1),
\]
y
\begin{equation}
 \boxed{
 \operatorname{Spec}\bigl(\operatorname{Sym}^2(F_{\rm av})\bigr)
 =\{\varphi^2,-1,\varphi^{-2}\}.}
 \label{p12-83-c9-s3:eq:espectro-sym2-phi}
\end{equation}
El único modo estable positivo es \(\varphi^{-2}\).
\end{theorem}

\begin{proof}
Los autovalores de una potencia simétrica de grado dos son los productos
\(\lambda_1^2\), \(\lambda_1\lambda_2\) y \(\lambda_2^2\). Con
\(\lambda_1=\varphi\) y \(\lambda_2=-\varphi^{-1}\), se obtienen
\(\varphi^2\), \(-1\) y \(\varphi^{-2}\). La primera excede uno, la
segunda tiene módulo uno y la tercera pertenece a \((0,1)\).
\end{proof}

Sobre una ruta cerrada con \(V_n=\varphi^n\), el modo estable satisface
\begin{equation}
 \boxed{M_n=\varphi^{-2n}=V_n^{-2}.}
 \label{p12-83-c9-s3:eq:modo-estable-phi}
\end{equation}

\subsection{Realización hiperbólica normalizada de la autoescala: \(\exp(2\log\varphi\,J_h)=F_{\mathrm{av}}^2\)}

Sea
\[
 A=F_{\rm av}^{\,2}=
 \begin{pmatrix}1&1\\1&2\end{pmatrix},
 \qquad
 J_h=\frac{2A-3I}{\sqrt5}.
\]
Entonces
\[
 J_h^2=I,
 \qquad
 A=\frac32I+\frac{\sqrt5}{2}J_h.
\]
Como
\[
 \cosh(2\log\varphi)=\frac32,
 \qquad
 \sinh(2\log\varphi)=\frac{\sqrt5}{2},
\]
se obtiene
\begin{equation}
 \boxed{
 \exp(2\log\varphi\,J_h)=F_{\rm av}^{\,2}.}
 \label{p12-83-c9-s3:eq:realizacion-hiperbolica-phi}
\end{equation}
La constante de autoescala es, por tanto, simultáneamente el autovalor de
Perron de la incidencia y el parámetro exponencial de su realización
hiperbólica normalizada.

\subsection{Caracterización, horquillas y prolongación}

La construcción finita determina el rayo positivo de \(F_{\rm av}\). La
caracterización analítica posterior es
\[
 x^2-x-1=0,
 \qquad x>0.
\]
Las horquillas \eqref{p12-83-c9-s3:eq:horquilla-fibonacci-phi} tienen extremos racionales
y anchura explícita \eqref{p12-83-c9-s3:eq:anchura-fibonacci-phi}. Para cada nivel
nonádico \(K\), elegimos el menor \(m=m_{\varphi}(K)\) cuya horquilla queda
contenida en un único subcilindro entre los \(729\) elementos de la fibra
ambiente. La minimalidad convierte la localización en una función, no en una
elección tácita.

Los cilindros seleccionados son compatibles con los truncamientos y sus
diámetros tienden a cero. Su intersección contiene un único real; el
\cref{p12-83-c9-s3:thm:rayo-perron-phi} lo identifica con \(\varphi\). Los cocientes de
Fibonacci son, pues, certificados racionales de la localización, no datos
decimales empleados para decidir la semilla.

\subsection{Certificado interno de unicidad, compatibilidad y convergencia de la autoescala}

\begin{criterion}[Refutación de la biografía de autoescala]
La construcción queda refutada si ocurre cualquiera de los hechos
siguientes:
\begin{enumerate}
 \item el predicado finito selecciona más de una órbita o más de una emisión;
 \item la recurrencia \(L_0,L_0,L_1,L_1\) no produce
       \eqref{p12-83-c9-s3:eq:w30-phi};
 \item el calendario no produce exactamente los tres pares usados en
       \eqref{p12-83-c9-s3:eq:N3-phi};
 \item se confunde un múltiplo de conteo \(N_{9},N_{27},N_{108}\) con una
       potencia de \(N_3\);
 \item \(F_{\rm av}\) posee otro rayo propio positivo;
 \item falla \eqref{p12-83-c9-s3:eq:potencias-Fav-fibonacci} o la identidad de Cassini;
 \item las horquillas racionales no son encajadas o no contraen su anchura;
 \item se identifica la frecuencia cronológica, la medida uniforme del
       catálogo o la medida de Parry;
 \item una tabla decimal interviene antes de construir el rayo de Perron;
 \item un cilindro seleccionado no trunca al cilindro anterior.
\end{enumerate}
\end{criterion}


% Unidad U015. Acoplamiento de horquillas racionales, refinamiento TPK y
% certificación a profundidad arbitraria.

\section{Certificación diagonal, límite inverso y generación a precisión arbitraria}
\label{p12-83-c9-s4:chap:precision-arbitraria-constantes}

\subsection{Separación tipada de los \(4\) índices de la prolongación nonádica}

Esta unidad tiene estatuto de certificación posterior. Los prefijos y sus
extensiones compatibles proceden de la relación interna
\(\operatorname{Ext}\), de la supervivencia coinductiva y de las leyes de
transición APP--TRIT--TPK. Las horquillas racionales que siguen no eligen
entre los \(729\) elementos de una fibra: localizan en la carta arquimediana
la extensión ya generada y certifican su precisión a profundidad arbitraria.

Para cada carácter \(\chi\in\{\pi,e,\varphi\}\) intervienen cuatro
parámetros de naturaleza diferente:
\begin{enumerate}
 \item \(n\), orden de la horquilla analítica racional;
 \item \(K\), nivel de refinamiento del Topological Phase Kernel;
 \item \(L_K=6K\), longitud ternaria del prefijo publicado;
 \item \(M\), número de posiciones decimales solicitado.
\end{enumerate}
No se postula \(n=K\), \(n=L_K\), \(K=M\) ni una proporcionalidad fija
entre ellos. El acoplamiento se construye mediante inclusiones de intervalos
decididas por aritmética entera.

Las partes enteras son
\[
 m_\pi=3,
 \qquad
 m_e=2,
 \qquad
 m_\varphi=1,
\]
y los caracteres fraccionarios son
\[
 \rho_\pi=\pi-3,
 \qquad
 \rho_e=e-2,
 \qquad
 \rho_\varphi=\varphi-1.
\]
Las tres familias de horquillas semiabiertas se escriben
\[
 J_{\chi,n}=[\ell_{\chi,n},h_{\chi,n}).
\]
Sus extremos son racionales, contienen a \(\rho_\chi\) y satisfacen
\[
 \ell_{\chi,n}\nearrow\rho_\chi,
 \qquad
 h_{\chi,n}\searrow\rho_\chi,
 \qquad
 h_{\chi,n}-\ell_{\chi,n}\longrightarrow0.
\]
Para \(\pi\), los extremos proceden de las sumas alternadas de la identidad
de Machin demostrada en el \cref{chap:biografia-pi-autonoma}; para \(e\),
de las sumas factoriales y su resto racional del
\cref{p12-83-c9-s2:chap:biografia-completa-e}; para \(\varphi\), de los cocientes de
Fibonacci y Cassini del \cref{p12-83-c9-s3:chap:biografia-completa-phi}.

\subsection{Cilindros ternarios y fibra ambiente inmediata}

Un prefijo ternario de longitud \(L_K=6K\) se representa por un entero
\(N_K\) con
\[
 0\leq N_K<3^{L_K}.
\]
Su cilindro semiabierto es
\begin{equation}
 C(N_K,L_K)=
 \left[
 \frac{N_K}{3^{L_K}},
 \frac{N_K+1}{3^{L_K}}
 \right).
 \label{p12-83-c9-s4:eq:cilindro-nivel-K}
\end{equation}
La siguiente emisión añade seis trits. La fibra ambiente inmediata contiene
\(3^6=729\) elementos, indexados por \(u\in\{0,\ldots,728\}\):
\begin{equation}
 C_{K,u}=
 \left[
 \frac{729N_K+u}{3^{L_K+6}},
 \frac{729N_K+u+1}{3^{L_K+6}}
 \right).
 \label{p12-83-c9-s4:eq:729-subcilindros}
\end{equation}
Los \(C_{K,u}\) son disjuntos y
\[
 C(N_K,L_K)=\bigsqcup_{u=0}^{728}C_{K,u}.
\]
El número \(729\) cuenta coordenadas de esta fibra ambiente; no cuenta
semillas, emisiones distintas, órbitas ni estados sobrevivientes.

\subsubsection{Comparación diofántica de los extremos racionales de la horquilla}

Sea \(J=[\ell,h)\subset C(N_K,L_K)\), con
\(\ell,h\in\mathbb Q\). Escribimos \(S_K=3^{L_K+6}\) y definimos
\begin{align}
 j_{\min}(J,K)
 &=\max\!\left(729N_K,\left\lfloor S_K\ell\right\rfloor\right),
 \label{p12-83-c9-s4:eq:jmin-diagonal}\\
 j_{\max}(J,K)
 &=\min\!\left(729N_K+728,\left\lceil S_Kh\right\rceil-1\right).
 \label{p12-83-c9-s4:eq:jmax-diagonal}
\end{align}

\begin{lemma}[Criterio de localización unitaria]
\label{p12-83-c9-s4:lem:criterio-localizacion-unitaria}
La horquilla \(J\) está contenida en un único elemento de la fibra
\eqref{p12-83-c9-s4:eq:729-subcilindros} si y sólo si
\begin{equation}
 \boxed{j_{\min}(J,K)=j_{\max}(J,K).}
 \label{p12-83-c9-s4:eq:jmin-igual-jmax}
\end{equation}
En ese caso, si \(j\) es el valor común,
\[
 u=j-729N_K,
 \qquad
 N_{K+1}=j.
\]
\end{lemma}

\begin{proof}
Los enteros \(j\) cuyos intervalos elementales de longitud \(S_K^{-1}\)
intersectan \(J\) forman el intervalo entero
\[
 \left[
 \lfloor S_K\ell\rfloor,
 \lceil S_Kh\rceil-1
 \right]\cap
 [729N_K,729N_K+728].
\]
Éste contiene exactamente un elemento si y sólo si sus extremos coinciden.
La igualdad identifica el numerador del subcilindro y, al truncar sus seis
últimos trits, se recupera \(N_K\).
\end{proof}

\subsection{Certificación diagonal mínima}

La generación no fija de antemano un mismo orden analítico para todos los
niveles. Una vez producida la extensión compatible, el orden de la
horquilla se aumenta sólo hasta certificarla de manera unitaria.

\begin{definition}[Recurrencia diagonal de certificación]
\label{p12-83-c9-s4:def:recurrencia-diagonal-constantes}
Fijada por \(\operatorname{Ext}\) una familia de cilindros compatibles
\(C(N_K,L_K)\) cuya publicación contiene a \(\rho_\chi\), se elige un orden
inicial \(n_\chi(K_0-1)\geq1\) y definimos recursivamente
\begin{equation}
 n_\chi(K)=
 \min\left\{
 n\geq n_\chi(K-1):
 j_{\min}(J_{\chi,n},K)=j_{\max}(J_{\chi,n},K)
 \right\},
 \label{p12-83-c9-s4:eq:nchi-K-minimo}
\end{equation}
y verificamos
\begin{equation}
 N_{\chi,K+1}
 =j_{\min}(J_{\chi,n_\chi(K)},K).
 \label{p12-83-c9-s4:eq:Nchi-K-mas-uno}
\end{equation}
El conjunto
\[
 \Delta_\chi=
 \{(n_\chi(K),K):K\geq K_0\}
\]
es la diagonal analítico--nonádica del carácter \(\chi\).
\end{definition}

\begin{theorem}[Existencia y unicidad de la diagonal de certificación]
\label{p12-83-c9-s4:thm:existencia-unicidad-diagonal}
Para \(\chi\in\{\pi,e,\varphi\}\), la recurrencia
\eqref{p12-83-c9-s4:eq:nchi-K-minimo} está definida en todo nivel finito. El índice
\(n_\chi(K)\) y la coordenada \(N_{\chi,K+1}\) son únicos.
\end{theorem}

\begin{proof}
Supongamos generada internamente la extensión cilíndrica de nivel \(K\) que contiene a
\(\rho_\chi\). Como \(\rho_\chi\) es irracional, no pertenece a ninguna
frontera ternaria \(3^{-L}\mathbb Z\). Por ello está en el interior de un
único subcilindro \(C_{K,u_*}\), a distancia positiva de sus dos extremos.
Las horquillas \(J_{\chi,n}\) contienen a \(\rho_\chi\) y su diámetro
tiende a cero; existe, por tanto, un \(n\) para el que la horquilla queda
contenida en \(C_{K,u_*}\). El conjunto del mínimo en
\eqref{p12-83-c9-s4:eq:nchi-K-minimo} es no vacío y, por el buen orden de
\(\mathbb N\), posee un menor elemento. El
\cref{p12-83-c9-s4:lem:criterio-localizacion-unitaria} prueba que la horquilla la
localiza de manera unitaria y que el numerador localizado coincide con el ya
generado. La inducción sobre \(K\) completa el argumento.
\end{proof}

\begin{corollary}[Ausencia de parámetros decimales]
La diagonal se calcula con sumas, productos, potencias, divisiones euclídeas,
pisos y techos de enteros. No recibe una cifra decimal de \(\pi\), \(e\) o
\(\varphi\) como dato de certificación.
\end{corollary}

\subsection{Compatibilidad con la conexión nonádica conjunta}

Sea \(\widehat x_{\chi,K}\) el estado enriquecido del canal \(\chi\) en
el nivel \(K\). Sus coordenadas incluyen el prefijo, el cilindro, la fase,
el cociente, el acarreo, la memoria vertical no acotada, el término terminal
y el registro de procedencia. La fase residual es
\[
 g_0(K)=[K]_9\in\mathbb Z/9\mathbb Z,
\]
mientras la etiqueta pública es
\[
 g(K)=1+((K-1)\bmod9)\in\{1,\ldots,9\}.
\]
La memoria vertical \(q_{\rm mem}(K)\in\mathbb Z_{\geq0}\) no se reduce
módulo \(12\); sólo su proyección dodecafásica
\[
 \beta(K)=q_{\rm mem}(K)\bmod12
\]
lo hace.

Para cada \(K\), sea
\[
 \Gamma_{9,K}:\widehat X_K\longrightarrow\widehat X_{K+9}
\]
la holonomía de una vuelta de fase. Entonces
\[
 g_0(K+9)=g_0(K),
 \qquad
 q_{\rm mem}(K+9)=q_{\rm mem}(K)+1,
 \qquad
 \beta(K+9)=\beta(K)+1\pmod{12}.
\]

\begin{theorem}[Retorno de fase con conservación genealógica]
\label{p12-83-c9-s4:thm:retorno-fase-conservacion-genealogica}
La igualdad de fase no reinicia el estado:
\[
 g_0(K+9)=g_0(K),
 \qquad
 \widehat x_{\chi,K+9}\ne\widehat x_{\chi,K}.
\]
Las profundidades \(27\), \(54\) y \(108\) son las composiciones tipadas
\begin{align*}
 \Gamma_{27,K}
 &=\Gamma_{9,K+18}\circ\Gamma_{9,K+9}\circ\Gamma_{9,K},\\
 \Gamma_{54,K}
 &=\Gamma_{9,K+45}\circ\cdots\circ\Gamma_{9,K},\\
 \Gamma_{108,K}
 &=\Gamma_{9,K+99}\circ\cdots\circ\Gamma_{9,K}.
\end{align*}
En ningún caso son reinicios del prefijo, del acarreo, del terminal o de la
memoria.
\end{theorem}

\begin{proof}
La primera coordenada se calcula en \(\mathbb Z/9\mathbb Z\), mientras
\(q_{\rm mem}\) pertenece a un monoide no acotado. Cada aplicación de
\(\Gamma_{9,K}\) añade una unidad de memoria, seis bloques por nivel
intermedio y la composición correspondiente de cociclos. Por ello la fase
puede coincidir aunque el estado enriquecido no lo haga. Las fórmulas de
profundidad muestran los índices de cada mapa y evitan componer como si todos
fueran endomorfismos del mismo espacio.
\end{proof}

\subsubsection{Compresión, expresión y término terminal}

La realización operatorial de una vuelta satisface
\begin{equation}
 U_\Gamma J=JT+\eta,
 \qquad
 J^*\eta=0,
 \qquad
 I-T^*T=\eta^*\eta.
 \label{p12-83-c9-s4:eq:compresion-expresion-precision}
\end{equation}
La contracción visible estricta se reserva para
\[
 M_9=\frac59V_9;
\]
no se identifica con el operador de compresión \(T\). La telescopía de
profundidad \(L\) conserva el término terminal:
\begin{equation}
 S_0=
 \sum_{r=0}^{L-1}M_9^{*r}D_9^*D_9M_9^r
 +M_9^{*L}S_0M_9^L.
 \label{p12-83-c9-s4:eq:telescopia-terminal-precision}
\end{equation}
Omitir el último sumando antes de probar su desaparición borraría
información de frontera y convertiría el retorno de fase en un reinicio
falso.

\subsection{Sistema inverso de historias compatibles}

Sea \(\mathcal H_K^\chi\) el conjunto de historias enriquecidas admisibles
del canal \(\chi\) a profundidad \(K\), y sea
\[
 p_{N,K}:\mathcal H_N^\chi\longrightarrow\mathcal H_K^\chi
 \qquad(N\geq K)
\]
el truncamiento. Definimos
\[
 \operatorname{Surv}_K^{(N)}(\chi)
 =p_{N,K}(\mathcal H_N^\chi)
\]
y
\begin{equation}
 \operatorname{Surv}_K(\chi)
 =\bigcap_{N\geq K}\operatorname{Surv}_K^{(N)}(\chi).
 \label{p12-83-c9-s4:eq:supervivencia-todas-profundidades}
\end{equation}
El objeto coinductivo es
\begin{equation}
 X_\chi=
 \varprojlim_K
 \bigl(\operatorname{Surv}_K(\chi),p_{K+1,K}\bigr).
 \label{p12-83-c9-s4:eq:limite-inverso-caracteres}
\end{equation}

\begin{theorem}[Sección compatible bajo prolongabilidad enriquecida]
\label{p12-83-c9-s4:thm:seccion-compatible-caracteres}
Supóngase además que, sobre cada subcilindro numérico generado, la
relación \(\operatorname{Ext}\) es no vacía, univaluada y compatible con los
truncamientos. Entonces \(\operatorname{Ext}\) y la supervivencia producen una familia
\[
 x_\chi=(\widehat x_{\chi,K})_{K\geq K_0}
\]
que satisface
\[
 p_{K+1,K}(\widehat x_{\chi,K+1})
 =\widehat x_{\chi,K}.
\]
Por tanto, \(x_\chi\in X_\chi\).
\end{theorem}

\begin{proof}
El subcilindro de nivel \(K+1\) tiene numerador
\(N_{\chi,K+1}=729N_{\chi,K}+u\). El truncamiento elimina sus seis últimos
trits y recupera \(N_{\chi,K}\). Las demás coordenadas se obtienen por las
leyes functoriales de acarreo, memoria, terminal y fase. La hipótesis de
univaluación de \(\operatorname{Ext}\) fija el subcilindro y la extensión
enriquecida; la unicidad local del
\cref{p12-83-c9-s4:lem:criterio-localizacion-unitaria} certifica después su
localización arquimediana, y la compatibilidad de \(\operatorname{Ext}\)
proporciona la igualdad de truncamiento.
\end{proof}

\subsection{Celdas decimales y precisión prescrita}

Para \(M\geq1\) y \(q\in\{0,\ldots,10^M-1\}\), sea
\[
 D_{M,q}=
 \left[
 \frac{q}{10^M},
 \frac{q+1}{10^M}
 \right).
\]
Una horquilla \(J=[\ell,h)\) está contenida en una única celda decimal si
y sólo si
\begin{equation}
 \boxed{
 \left\lfloor10^M\ell\right\rfloor
 =\left\lceil10^Mh\right\rceil-1=q.}
 \label{p12-83-c9-s4:eq:criterio-celda-decimal}
\end{equation}

\begin{definition}[Nivel mínimo de certificación decimal]
Para \(\chi\in\{\pi,e,\varphi\}\), definimos \(K_\chi(M)\) como el menor
nivel para el que existen un orden analítico \(n\) y un entero \(q\) tales
que
\begin{enumerate}
 \item \(J_{\chi,n}\) satisface
       \eqref{p12-83-c9-s4:eq:criterio-celda-decimal};
 \item \(J_{\chi,n}\) queda contenido en el cilindro generado por el
       estado \(\widehat x_{\chi,K}\);
 \item la transición al nivel siguiente satisface
       \eqref{p12-83-c9-s4:eq:jmin-igual-jmax}.
\end{enumerate}
El par \((K,n)\) se minimiza lexicográficamente.
\end{definition}

\begin{theorem}[Certificación a precisión decimal arbitraria de las salidas coinductivas]
\label{p12-83-c9-s4:thm:precision-decimal-arbitraria}
Para cada \(\chi\in\{\pi,e,\varphi\}\) y cada \(M\geq1\), existe
\(K_\chi(M)<\infty\). La sección generada determina un único prefijo decimal
de \(M\) posiciones, y la horquilla reconoce y certifica ese prefijo. Si
\(M'<M\), el prefijo de orden \(M\) trunca al de
orden \(M'\), y las celdas correspondientes están encajadas.
\end{theorem}

\begin{proof}
Los tres caracteres son irracionales: \(\varphi\) lo es por su polinomio
cuadrático con discriminante no cuadrado; la irracionalidad de \(e\) y
\(\pi\) se toma de sus teoremas clásicos de caracterización, posteriores a
la construcción finita. Ninguno pertenece, pues, a una frontera decimal o
ternaria racional. Para un \(M\) fijo, la distancia del carácter al conjunto
finito de fronteras decimales relevantes es positiva. Las horquillas
racionales contraen su diámetro hasta quedar en una única celda decimal.
El \cref{p12-83-c9-s4:thm:existencia-unicidad-diagonal} acopla esa horquilla con un
cilindro TPK y con una extensión única. El mínimo lexicográfico existe por
buen orden. La compatibilidad de la sección prueba el encajamiento al
aumentar \(M\).
\end{proof}

\begin{corollary}[Unicidad real]
Las celdas encajadas tienen diámetros \(10^{-M}\to0\). Su intersección es
un singleton, identificado por las caracterizaciones precedentes con
\(\pi\), \(e\) o \(\varphi\), respectivamente.
\end{corollary}

\subsection{Equivalencia de \(2\) realizaciones finitas de la ley uniforme}

Los testigos de \(1,000\) y \(10,000\) posiciones no definen los
caracteres ni limitan la prolongación. Son ejecuciones de la misma recurrencia
uniforme con dos profundidades distintas.

\begin{table}[htbp]
\centering
\caption{Dimensiones exactas de las ejecuciones certificadas, por canal.}
\label{p12-83-c9-s4:tab:dimensiones-ejecuciones-1000-10000}
\begin{tabular}{lrr}
\toprule
Objeto contado & \(1,000\) posiciones & \(10,000\) posiciones\\
\midrule
Particiones exhaustivas de la fibra ambiente & 396 & 3501\\
Trits transportados & 2376 & 21006\\
Pares de igual fase \((K,K+9)\) & 387 & 3492\\
Vueltas nonádicas completas & 43 & 388\\
Elementos examinados por partición & 729 & 729\\
Extensiones compatibles conservadas por partición & 1 & 1\\
\bottomrule
\end{tabular}
\end{table}

Para \(10,000\) posiciones,
\[
 B=3501=389\cdot9,
 \qquad
 6B=21006,
 \qquad
 3^{6B}>10^{10000}.
\]
Los primeros \(396\) bloques de la ejecución larga coinciden con la
ejecución de \(1,000\) posiciones. En cada uno de los \(3492\) pares de
igual fase se preserva la coordenada de fase y se prolongan el prefijo, la
profundidad, el cilindro y la memoria. El certificado registra también las
repeticiones eventuales de proyecciones locales sin confundirlas con retorno
del estado completo.

\subsection{Separación entre generación y verificación}

\begin{definition}[Generador estructural]
El generador recibe el catálogo finito, \(L_0,L_1\), el estado inicial, las
leyes de conexión y la relación interna \(\operatorname{Ext}\). En cada nivel:
\begin{enumerate}
 \item enumera los \(729\) elementos de la fibra ambiente;
 \item aplica \(\operatorname{Ext}\) y conserva la única extensión compatible;
 \item actualiza fase, cociente, acarreo, memoria y terminal;
 \item escribe el bloque ternario, el recibo de selección y el registro genealógico.
\end{enumerate}
\end{definition}

\begin{definition}[Comprobación posterior independiente]
La comprobación posterior independiente recibe las salidas cerradas del
generador y las tres familias de horquillas racionales; calcula
\(j_{\min}\) y \(j_{\max}\), certifica su igualdad con la extensión ya
producida, reconstruye las
particiones, comprueba los cuadrados de truncamiento, los pares
\((K,K+9)\), las sumas de partición
\[
 \operatorname{killed}_{\mathrm{izq}}+1+\operatorname{killed}_{\mathrm{der}}=729,
\]
la integridad de las salidas y, al final, la coincidencia con expansiones de referencia.
\end{definition}

\begin{theorem}[Aislamiento de los datos objetivo]
\label{p12-83-c9-s4:thm:aislamiento-datos-objetivo-precision}
Las expansiones decimales de referencia y los datos metrológicos no
intervienen en la selección de una extensión compatible. La comparación se
ejecuta después de escribir los tres prefijos.
\end{theorem}

\begin{proof}
La entrada matemática de cada decisión generativa consiste en el catálogo
finito, \(L_0,L_1\), el estado heredado y la relación interna
\(\operatorname{Ext}\). Los extremos racionales, potencias de \(3\), pisos y
techos pertenecen a la comprobación posterior. El generador termina antes de
que ésta abra las horquillas o ejecute los controles decimales. La dirección
del flujo de datos es
\[
 \text{estructura}\longrightarrow
 \text{prefijos y ledgers}\longrightarrow
 \text{comprobación posterior independiente};
\]
no existe una flecha en sentido contrario.
\end{proof}

\subsection{Obstrucciones diferenciadas de la generación y de su certificación a precisión arbitraria}

\begin{criterion}[Refutación de la generación a precisión arbitraria]
Los fallos de órbita, extensión, truncamiento o retorno refutan la generación;
los fallos de horquilla, orden o control decimal refutan el certificado
analítico correspondiente; y una discrepancia entre ambos refuta la
identificación conjunta de la salida publicada. Con esta separación de tipos,
la afirmación completa queda refutada si:
\begin{enumerate}
 \item se identifica cualquiera de los índices \(n,K,L_K,M\) sin un
       teorema de enlace;
 \item una horquilla deja de contener su carácter o no contrae su diámetro;
 \item una etapa satisface \(j_{\min}\ne j_{\max}\) después del orden
       declarado como mínimo;
 \item una extensión no trunca al cilindro del nivel anterior;
 \item el retorno de fase reinicia prefijo, acarreo, terminal o memoria;
 \item se omite el término terminal en
       \eqref{p12-83-c9-s4:eq:telescopia-terminal-precision} sin demostrar su anulación;
 \item los primeros \(396\) bloques de la ejecución larga difieren de la
       ejecución corta;
 \item un control decimal o metrológico es leído antes de cerrar la salida;
 \item una cifra impresa difiere de la sección matemática certificada;
 \item la prueba depende de que la precisión solicitada sea \(1,000\) o
       \(10,000\), en vez de un \(M\) arbitrario.
\end{enumerate}
\end{criterion}
